\documentclass[10pt,a4paper]{amsart}
\usepackage[margin=19mm]{geometry}
\usepackage{amsmath,amssymb,mathtools}
\usepackage[scr=rsfs,cal=euler]{mathalfa}
\usepackage{xfrac}
\usepackage[unicode,hidelinks,bookmarks=false]{hyperref}
\newcommand{\ie}{i.e.\ }
\newcommand{\cf}{cf.\ }
\newcommand{\resp}{respectively\ }
\newcommand{\Subsection}{\subsection}
\providecommand{\nobreakdash}{\nobreak\hbox{-}}
\setlength{\emergencystretch}{2em}
\multlinegap0pt
\usepackage{bm}
\usepackage{bbm}
\usepackage{dsfont}
\usepackage{xspace}
\usepackage{cancel}
\usepackage{mathtools}
\usepackage{cleveref}
\usepackage{amsthm}
\makeatletter
\@ifundefined{theorem}{\newtheorem{theorem}{Theorem}}{}
\@ifundefined{lemma}{\newtheorem{lemma}[theorem]{Lemma}}{}
\@ifundefined{proposition}{\newtheorem{proposition}[theorem]{Proposition}}{}
\makeatother
\makeatletter
\def\I{{ \mathcal{I}}}
\def\L{{ \mathcal{L}}}
\def\O{{ \mathcal{O}}}
\def\QQ{\overline{\mathbb Q}}
\def\R{{\mathbb R}}
\def\Z{{\mathbb Z}}
\expandafter\ifx\csname claimproof\endcsname\relax
\newenvironment{claimproof}[0]
  {%
   \paragraph{\it Proof.}%
  }
  {%
    \hfill$\blacksquare$%
  }
\fi
\def\dd{{ \mathfrak{d}}}     
\newcommand{\e}{\varepsilon}
\newcommand{\g}{\gamma}
\makeatother
\title[Growth of generalized greatest common divisors along orbits ]{Growth of generalized greatest common divisors along orbits of self-rational maps on projective varieties}
\author{Yohsuke Matsuzawa}
\date{Source: arXiv:2507.05027v1 (CC BY 4.0)}
\begin{document}
\maketitle

\noindent\emph{Source and licence.} Growth of generalized greatest common divisors along orbits of self-rational maps on projective varieties. Version arXiv:2507.05027v1 (CC BY 4.0), distributed under the Creative Commons Attribution 4.0 International licence, as recorded in the arXiv licence metadata for this exact version and shown on the abstract page https://arxiv.org/abs/2507.05027v1. Full rights notice: licenses/arxiv-2507.05027-CC-BY-4.0.txt.\par\smallskip

\noindent\textit{Editorial scope.} Counted unit: the Proposition whose source label is \texttt{prop:ggcdpullbackbound} in the article (source0.tex lines 1688--1701, source label \texttt{prop:ggcdpullbackbound}), delivered together with the article's own complete original proof of it (source0.tex lines 1702--1758, one \texttt{proof} environment of 57 lines). No mathematics is written, repaired or re-derived: statement and proof are the article's own text line for line. Required context retained uncounted: lem:glsecthick (lines 752--767); prop:existencevansectau (lines 954--974); eq:condonrm (lines 967--969); prop:boundoftaudyn (lines 1274--1287). Every definition, display and label of those lines is retained unchanged. Omitted from this package and not redistributed: 1--751, 768--953, 970--1273, 1288--1687, 1759--2056. Nothing inside a retained excerpt is omitted or altered: every retained body is delivered exactly as the article's lines are written, so no omission receipt of any kind accompanies this package. Citations: the retained excerpts carry no citation command at all, so no bibliography entry of the article is transcribed and no reference is redistributed. Reference adaptation: none was needed and none was made; every reference inside the delivered unit resolves inside it, so no unresolved reference or citation is produced. Printed slips kept literal: no printed word, symbol, hypothesis or number of the counted statement or of its proof was altered. Layout and class: the article's own class is \texttt{amsart} with packages including etex, ifpdf, geometry, iftex, luatex85, amsaddr, cite, tikz, tikz-cd, lipsum, adjustbox, fancybox, listings, docmute; the wrapper is the lane's pinned \texttt{amsart} editorial wrapper at 10pt on A4, which restates the theorem-like environments the retained text needs and adds the article's own macro definitions that the retained text needs (\texttt{\textbackslash I}, \texttt{\textbackslash L}, \texttt{\textbackslash O}, \texttt{\textbackslash QQ}, \texttt{\textbackslash R}, \texttt{\textbackslash Z}, \texttt{\textbackslash dd}, \texttt{\textbackslash e}, \texttt{\textbackslash g}), with extra wrapper packages (bm, bbm, dsfont, xspace, cancel, mathtools, cleveref). The delivered numbering is the unit's own. Assets: no retained span contains a figure environment, a graphics-inclusion command, a PDF-page inclusion, a table or a TikZ picture; no image file is copied into this package, the package declares no asset dependency and no image file is redistributed with it. Licence: version arXiv:2507.05027v1 of this article is distributed under the Creative Commons Attribution 4.0 International licence, as recorded in the arXiv licence metadata for this exact version and shown on the abstract page https://arxiv.org/abs/2507.05027v1. A full rights notice is delivered at licenses/arxiv-2507.05027-CC-BY-4.0.txt. Characters: every retained excerpt is pure ASCII, so the delivered engine, XeLaTeX with its default Latin Modern text font, reports no missing character.
\begin{lemma}\label{lem:glsecthick}
Let $k$ be an infinite field.
Let $X$ be a projective variety over $k$ of dimension $N \geq 0$.
Let $Y \subset X$ be a proper closed subscheme of pure dimension $l$ with ideal sheaf $ \I \subset \O_{X}$.
Let $ \mathcal{L}$ be a very ample invertible $\O_{X}$-module on $X$
and $\dd \subset |\L|$ be a linear system that induces a closed immersion.
Then for a general sequence $H_{1}, \dots, H_{l} \in \dd$, we have
$H_{1\cdots i}:= H_{1} \cap \dots \cap H_{i}$ is irreducible and reduced of dimension $N-i$,
$Y \cap H_{1\cdots i}$ is pure dimension $l - i$, and
\begin{align}
h^{0}(X, \mathcal{L}^{\otimes m} \otimes \O_{X}/\I^{r}) \leq \sum_{i=0}^{l} \binom{m+i-1}{i} h^{0}(H_{1\cdots i}, \O_{H_{1\cdots i}}/\I_{i}^{r})
\end{align}
for all $m, r \in \Z_{\geq 1}$,
where $\I_{i} \subset \O_{H_{1\cdots i}}$ is the ideal sheaf of $Y \cap H_{1\cdots i} \subset H_{1\cdots i}$.
Note that $H_{1\cdots i} = X$ when $i=0$ by convention.
\end{lemma}

\begin{proposition}\label{prop:existencevansectau}
Let $k$ be an algebraically closed field of characteristic zero.
Let $X$ be a projective variety over $k$ of dimension $N \geq 0$.
Let $Y \subset X$ be a proper closed subscheme of pure dimension $l$ with ideal sheaf $ \I \subset \O_{X}$.
Let $ \mathcal{L}$ be a very ample invertible $\O_{X}$-module on $X$.
Consider a general sequence $H_{1}, \dots, H_{l} \in | \mathcal{L}|$ such that the conclusions of \cref{lem:glsecthick} hold.
Set 
\begin{align}
\tau_{i} = \tau(Y\cap  H_{1\cdots i} ,  H_{1\cdots i})
\end{align}
for $i=0,\dots, l$.
Let $\e_{0} > 0$.
Then there is $a \in \Z_{\geq 1}$ such that for all $m,r \geq a$ with
\begin{align}
\sum_{i=0}^{l} \frac{\tau_{i}}{i !} \bigg( \frac{r}{m} \bigg)^{N-i} \leq \frac{ (c_{1}( \mathcal{L})^{N}) }{N!} - \e_{0}, \label{eq:condonrm}
\end{align}
we have
\begin{align}
H^{0}(X, \mathcal{L}^{\otimes m} \otimes \I^{r}) \neq 0.
\end{align}
\end{proposition}

\begin{proposition}\label{prop:boundoftaudyn}
Let $k$ be an algebraically closed field of characteristic zero.
Let $X$ be a regular projective variety over $k$ of dimension $N \geq 0$.
Let $f \colon X \longrightarrow X$ be a surjective morphism.
Let $Y \subset X$ be a proper closed subscheme of pure dimension $l$.
Suppose $Y \subset X$ is a regular embedding.
Let $H$ be a very ample divisor on $X$.
Then there is $C > 0$ depending only on $X, Y, H$ such that
for any $n \geq 0$, $i=0,\dots, l$, and a general sequence $H_{1}, \dots, H_{i} \in |H|$
(general subject to $n$), we have
\begin{align}
\tau(f^{-n}(Y) \cap H_{1} \cap \cdots \cap H_{i}, H_{1} \cap \cdots \cap H_{i}) \leq C ((f^{n})^{*} H^{N-i} \cdot H^{i}  ).
\end{align}
\end{proposition}

\begin{proposition}\label{prop:ggcdpullbackbound}
Let $X$ be a smooth projective variety over $\QQ$ of dimension $N$.
Let $f \colon X \dashrightarrow X$ be a dominant rational map.
Let $Y \subset X$ be a proper closed subscheme of pure dimension $l$.
Suppose $Y \subset X$ is a regular embedding and $Y \subset X_{f}^{\rm back}$.
Let $H$ be a very ample divisor on $X$.
Then for any $\e >0$, there is $n_{0} \in \Z_{\geq 0}$
such that for all $n \geq n_{0}$, there is a proper closed subset $Z \subset X$
and $\g \in \R$ such that
\begin{align}
h_{f^{-n}(Y)} \leq (d_{N-l}(f)^{1/(N-l)} + \e)^{n} h_{H} + \g
\end{align}
on $(X\setminus Z)(\QQ) $.
\end{proposition}

\begin{proof}
Let us fix arbitrary $\e>0$.
By log concavity of dynamical degrees, we have 
$d_{N-i}(f)^{1/(N-i)} \leq d_{N-l}^{1/(N-l)}$ for $i = 0,\dots, l$.
Thus there is $\e'>0$ such that
\begin{align}
( d_{N-i}(f)+ \e' )^{1/(N-i)} <  d_{N-l}(f)^{1/(N-l)}+ \frac{\e}{2} \label{eq:dyndegineq}
\end{align}
for $i=0,\dots,l$.

Let $C > 0$ be as in \cref{prop:boundoftaudyn}.
Then there is $n_{0} \in \Z_{\geq 0}$ such that for all $n \geq n_{0}$ and $i=0,\dots,l$, we have
\begin{align}
C ((f^{n})^{*}(H^{N-i}) \cdot H^{i} )  \leq ( d_{N-i}(f)+ \e')^{n}.
\end{align}

Let $\e_{0} > 0$ be such that $(H^{N})/N! - \e_{0} > 0$. 
Recalling  \cref{eq:condonrm}, we consider the following condition
\begin{align}
\sum_{i=0}^{l} \frac{( d_{N-i}(f)+ \e')^{n}}{i!}  \frac{1}{( d_{N-l}(f)^{1/(N-l)}+ \e/2 )^{n(N-i)} } \leq \frac{(H^{N})}{N!} - \e_{0}.
\end{align}
By \cref{eq:dyndegineq},  enlarging $n_{0}$ if necessary, this inequality holds for all $n \geq n_{0}$.

Now pick arbitrary $n \geq n_{0}$.
We apply \cref{prop:existencevansectau} to $X$, $f^{-n}(Y)$, $\O_{X}(H)$, and $\e_{0}$.
Let $H_{1}, \dots, H_{l} \in |H|$ be a general sequence such that the conclusions of
\cref{prop:existencevansectau}  and \cref{prop:boundoftaudyn} hold.
Let $a$ be the one obtained by \cref{prop:existencevansectau}.
Then for any integers $r,m \geq a$ such that
\begin{align}
\frac{r}{m} \leq \frac{1}{( d_{N-l}(f)^{1/(N-l)}+ \e/2 )^{n} }, \label{eq:rmadmrange}
\end{align}
we have
\begin{align}
&\sum_{i=0}^{l} \frac{\tau(f^{-n}(Y)\cap H_{1\cdots i}, H_{1\cdots i}  )}{i!} \bigg( \frac{r}{m} \bigg)^{N-i}\\
&\leq 
\sum_{i=0}^{l} \frac{( d_{N-i}(f)+ \e')^{n}}{i!}  \frac{1}{( d_{N-l}(f)^{1/(N-l)}+ \e/2 )^{n(N-i)} } \leq \frac{(H^{N})}{N!} - \e_{0},
\end{align}
where $H_{1\cdots i} = H_{1} \cap \cdots \cap H_{i}$.
Thus, if $\I \subset \O_{X}$ is the ideal sheaf of $f^{-n}(Y) \subset X$, we have
\begin{align}
H^{0}(X, \O_{X}(mH) \otimes \I^{r}) \neq 0
\end{align}
whenever $r, m \geq a$ and they satisfy \cref{eq:rmadmrange}.
Therefore there is an effective divisor $D$ linearly equivalent to $mH$ containing 
the closed subscheme defined by $\I^{r}$.
This implies 
\begin{align}
rh_{f^{-n}(Y)} \leq mh_{H} + O(1)
\end{align}
on $(X \setminus D)(\QQ)$.
Taking $r,m \geq a$ so that 
\begin{align}
\frac{m}{r} \leq (d_{N-l}(f)^{1/(N-l)} + \e)^{n},
\end{align} 
we are done.
\end{proof}

\end{document}
